\subsection{Size and accuracy}
\label{sec:graphres}
Across the \GraphTextOK{} judgments the extractor found \GraphUnits{} distinct pairs of a citing judgment and a reporter citation of another judgment. \GraphResolved{} of them (\GraphResolvedPct) were resolved: \GraphExact{} by number, \GraphParallel{} through a parallel number and \GraphName{} by case name and year. \GraphConflict{} pairs whose number and name disagreed were dropped. The resolved pairs form \GraphEdges{} distinct links between \GraphCiting{} citing and \GraphCited{} cited judgments, so \GraphCitedPct{} of the Supreme Court's reported judgments are cited by at least one later reported judgment. Resolution depends on the era of the citing judgment (Table~\ref{tab:graph}). Until the 1980s judgments cite mostly the S.C.R. and are resolved by number. From the 1990s the commercial SCC is the most cited report. Parallel S.C.R. citations, which make recent judgments easy to resolve, become common only in the 2010s, so the 2000s are the hardest decade.

\begin{table}[t]
\caption{Citations of other reported judgments, by decade of the citing judgment, and how they were resolved.}
\label{tab:graph}
\centering
\footnotesize
\setlength{\tabcolsep}{4pt}
\begin{tabular}{lrrrrr}
\toprule
Citing & Citations & Resolved & \multicolumn{3}{c}{Resolved by} \\
\cmidrule(lr){4-6}
decade & & (\%) & Number & Parallel & Name \\
\midrule
1950s & 953 & 87.6 & 821 & 0 & 14 \\
1960s & 3{,}259 & 84.0 & 2{,}669 & 6 & 63 \\
1970s & 5{,}001 & 75.9 & 3{,}561 & 18 & 219 \\
1980s & 6{,}467 & 77.5 & 4{,}160 & 43 & 806 \\
1990s & 15{,}667 & 63.8 & 4{,}958 & 361 & 4{,}676 \\
2000s & 34{,}810 & 55.0 & 2{,}957 & 375 & 15{,}830 \\
2010s & 70{,}340 & 78.5 & 23{,}027 & 17{,}174 & 15{,}026 \\
2020s & 67{,}418 & 84.0 & 26{,}370 & 22{,}036 & 8{,}212 \\
\bottomrule
\end{tabular}

\end{table}

\textbf{Accuracy.} Where a citation carries an S.C.R. or neutral number and the case name beside it can also be resolved, the name and year alone reach the same record as the number in \GraphNameSamePct{} of \GraphNamePrecN{} cases. In another \GraphNameLoosePct{} they reach a different record with matching parties (for example another order in the same case), and in \GraphNameWrongPct{} a record with different parties. This is a label-free estimate of the precision of name-only links. Only \GraphLater{} links (\GraphLaterPct) point to a judgment decided after the citing one, which is impossible. In a sample of them the causes were wrong decision dates in the metadata, page numbers of a neighbouring judgment in the same report volume, and truncated generic names (\emph{Others v.\ The State of Assam}) matched to the wrong judgment. Because name matching is restricted to the cited year, most wrong links cannot point forward, so this check is a sanity check, not an error bound. We also checked by hand \GraphCheckN{} randomly drawn name-only links from judgments outside the pilot, against the citation context and the index record: all \GraphCheckOK{} were correct, including names misspelt in the citing text, such as \emph{A.N.\ Shehgal} for \emph{A.N.\ Sehgal}. Of \UnresN{} randomly drawn SCC and AIR citations that stayed unresolved, \UnresNotIndex{} cite judgments that are not in the index at all, because they were reported in the SCC or AIR but not in the S.C.R. The other \UnresInIndex{} could have been resolved: \UnresExtraction{} were lost to OCR errors or truncated names in the citing text, \UnresSource{} to errors in the source (a wrong year in the citation, a misspelt index title), \UnresAmbig{} to names shared by several judgments, and \UnresVerifier{} to limitations of the frozen name parser.

\textbf{Face validity.} The most cited judgments are \emph{Maneka Gandhi v.\ Union of India} (\GraphTopOneN{} citing judgments), \emph{Kesavananda Bharati} (\GraphTopTwoN), \emph{Sharad Birdhi Chand Sarda v.\ State of Maharashtra} (\GraphTopThreeN), \emph{Ramana Dayaram Shetty v.\ International Airport Authority} (\GraphTopFourN) and \emph{Shivaji Sahebrao Bobade v.\ State of Maharashtra} (\GraphTopFiveN): leading authorities on personal liberty, the basic structure, circumstantial evidence, State action and criminal appeals. Of the \GraphLandmarkN{} landmark judgments that we chose by hand for the benchmark before the graph existed, \GraphLandmarkCited{} are cited at least once, and their median in-degree is higher than that of \GraphLandmarkPct\% of all judgments.
